\begin{frame}{마지막 체크리스트 — 학기 이후에 가져갈 세 가지}
  {\footnotesize 공식 목록은 13개 장이지만, 졸업하고 몇 년 뒤에
  실제로 남는 것은 \textbf{세 가지}:}
  \vskip0.4em
  \begin{enumerate}
    \item \kb{``무엇을 데이터(경험)으로 갖고, 무엇을 최적화하고
      싶은가''라는 질문의 습관.} 새로운 문제(운영 전략, A/B 테스트,
      물리 설계, LLM 후학습)를 마주했을 때 ``어떤 에이전트·어떤
      보상·어떤 경험인가''로 \textbf{번역하는 연습}을 13번 한 것이
      이 학기. 알고리즘 이름은 잊어져도 \textbf{이 번역은 남는다} —
      ML1의 ``행의 단위, 정답의 유무''와 쌍을 이루는, 문제
      정식화 기술의 절반이 완성된 것.
    \item \kb{정직한 숫자에 대한 기준.} 시드 $\geq 3$ · 학습/평가
      분리 · 베이스라인 · ``완전히 수렴하지 않았다''는 문장 — 8.1
      에서 시작해 16.1에서 반복된 프로토콜의 정직성은, 강화학습만이
      아니라 \textbf{모든 실험 과학}(ML1 val/test 정직성과 같은
      뿌리)의 기본 문법.
    \item \kb{``같은 방정식, 다른 형태''의 지도.} 이 학기의 알고리즘
      13개는 \textbf{3단계 구조의 13가지 변주}에 불과 — 그래서 새
      알고리즘 논문을 만나도 ``이것은 ②번(수치화)을 바꾼 변주구나''
      에서부터 읽을 수 있다. 지도의 \textbf{엣지}를 말할 수 있는
      한, 지도는 \textbf{확장 가능한 자산}이고, \textbf{박스만}
      외운 한으로는 \textbf{재고}에 불과하다.
  \end{enumerate}
\end{frame}
